\section{Experimental Setup}
\label{app:setup}

\paragraph{Environments.}
\emph{Toy POMDP:} $|S|{=}2$, $|A|{=}2$, $|O|{=}3$, $T{=}3$, with closed-form
oracle bridge functions obtained by pseudo-inverse and exact dynamic-programming
policy values. This is the only substrate where recovery error against ground
truth is directly measurable, and it carries every machine-precision correctness
check.

\emph{Benchmark:} 720 observable-core states $\times$ 2 latent states, 8 actions,
$T{=}3$, built on the public \texttt{gumbel-max-scm} clinical simulator
\citep{oberst2019}. The latent index drives dynamics and behavior but is censored
from the logs. We construct a pre-decision negative control $O_0$ satisfying the
identification assumption by construction, plus a per-step noisy marker that
guarantees bridge existence --- deterministic masking alone gives rank
$720 < 1440$, i.e.\ no valid one-step bridge. A confounding knob
$\kappa \in \{0, 0.25, 0.5, 0.75, 1.0\}$ sweeps from no confounding to strong
confounding.

\emph{Continuous:} scalar linear-Gaussian confounded POMDP, described in
\cref{app:cont}.

\paragraph{Protocol.} All multi-seed results use seeds $\{0,1,2,3,4\}$ and report
$1.96\,\hat\sigma/\sqrt{n}$ intervals. Ground truth throughout is exact dynamic
programming on the latent SCM (tabular) or a closed-form recursion (continuous).
The full pipeline runs on a laptop in NumPy; no GPU is used or required.

\paragraph{Scaling the benchmark estimator.} A direct implementation of the
paper's history-conditioning would require on the order of billions of table
cells at benchmark scale. Three exact structural factorizations specific to the
environment's observation design --- hashed encoding, core-diagonal bridge
support, and closed-form $2\times2$ local inversions --- keep the pipeline
laptop-scale without approximation.

\section{Continuous Extension: Derivation and Details}
\label{app:cont}

\subsection{Environment}
State, action, observation and reward are all real-valued. The observation is a
noisy proxy $o_t = s_t + \varepsilon_t$ for the hidden state, and the logging
policy may act on the true $s_t$ (confounded) or only on $o_t$ (not confounded),
mixed by a knob $\kappa$. This is the classical \emph{measurement-error} setting
that proximal causal inference was designed for: fitting a model that treats the
noisy proxy as if it were the state induces errors-in-variables bias whenever
$\kappa > 0$.

Because the reward is linear and candidate policies are affine and
observation-only ($a = K o + c$), the closed-loop mean trajectory is an exact
scalar recursion --- no covariance propagation is needed. This yields an exact
closed-form oracle for both policy values and bridge coefficients, verified to
reproduce each other to $1.665\times10^{-16}$.

\subsection{Mean-value bridges}
The paper's bridge functions are \emph{distributional}: they must match
$p(y \mid X_t)$ for every $y$. Reproducing that in continuous $y$ requires
conditional density estimation. Since the downstream quantity is a policy-value
\emph{mean} and the pessimism machinery concerns estimation uncertainty in the
bridge coefficients rather than the environment's intrinsic reward noise, a
conditional-mean moment restriction suffices and is exact for that purpose:
\begin{equation}
\mathbb{E}\!\left[b(A_t, O_t) \mid X_t = x\right]
= \mathbb{E}\!\left[Y_t \mid X_t = x\right] \quad \forall x,
\end{equation}
with $Y_t = R_t$ (reward bridge) or $O_{t+1}$ (dynamics bridge). This is exactly
a nonparametric instrumental-variable regression, with $W = (A, O)$ confounded,
$X = (A, H_{t-1}, O_0)$ the instrument, and $Y$ the outcome --- well precedented
in the proximal literature \citep{mastouri2021}. It is a deliberate scope
reduction, not a fundamental barrier.

\subsection{Two-stage kernel solve}
We derived the closed form directly from a stated objective rather than
pattern-matching a remembered formula, because a mis-remembered index convention
would silently produce a wrong-but-plausible estimator --- the exact failure mode
oracle verification exists to catch. With $K_X$, $K_W$ the instrument- and
bridge-side kernels,
\begin{align}
\Gamma &= \left(K_X(X_1,X_1) + N_1\lambda_1 I\right)^{-1} K_X(X_1,X_2), \\
Z &= K_W(W_1,W_1)\,\Gamma, \\
H &= \tfrac{1}{N_2} Z Z^\top + \lambda_2 K_W(W_1,W_1), \quad
g = \tfrac{1}{N_2} Z y_2,
\end{align}
and $\beta = H^{-1} g$. \textbf{Unifying check:} when $K_W$ is a one-hot/delta
kernel (the tabular case), $K_W(W_1,W_1) = I$ exactly and this reduces to the
tabular estimator's own primal solve --- a correctness check on the derivation
itself, not merely an analogy.

\subsection{A rank-deficiency pitfall}
For the primary verified path we take $K_W$ to be a \emph{linear} kernel, since
the true bridges are exactly linear and a linear kernel's RKHS therefore contains
them. But a linear kernel has a finite ($2$-dimensional) feature space, so the
$N_1 \times N_1$ Gram matrix $K_W(W_1,W_1)$ has rank $\le 2$ --- confirmed
empirically, with the remaining eigenvalues at the $\sim\!10^{-13}$ floor. The
dual system is therefore numerically meaningless in those directions.

Notably, the specific projected quantity we originally extracted
($\hat\theta = W_1^\top \beta$) turned out to be \emph{invariant} to that
instability and produced identical recovery numbers before and after the fix ---
so the bug was invisible from the recovery metrics. It still had to be fixed,
because $H$ and $\hat b$ enter the pessimism ellipsoid directly, where no such
invariance protects them. The fix is to solve in the $2$-dimensional
\emph{primal} coordinates: mathematically identical, exactly well-conditioned,
and incidentally $\sim\!1000\times$ cheaper.

\subsection{Spectral monitoring}
The tabular fix for a misleading condition number was a hard eigengap on the
smallest signal eigenvalue, correct there because the null space is
\emph{structural} ($|O| > |S|$, exactly zero signal beyond it). A continuous
kernel Gram matrix has smooth spectral decay rather than a structural wall, so a
hard eigengap does not transfer. We use the standard kernel-ridge
\emph{effective dimension} $N_{\mathrm{eff}}(\lambda) = \mathrm{tr}\,K(K + N\lambda I)^{-1}$
\citep{caponnetto2007} instead.

\subsection{Gradients}
The value chain is multilinear in each stage's coefficients, so reward-block
gradients are exact. A dynamics block's coefficients additionally shift the
\emph{evaluation point} of every later stage, and that dependence runs through
the kernel nonlinearly; we keep the coefficient-linear part exact and obtain the
one remaining scalar sensitivity by a single 1-D central difference --- $O(T)$
extra work per gradient call, not $O(N_1)$ --- verified against numerical
differentiation of the whole chain to $\approx\!10^{-11}$.

\subsection{Regularization: a tested-and-rejected hypothesis}
We hypothesized that the tabular-calibrated $\lambda_1 = 1/N_1$ would be too weak
for a genuine kernel conditional mean embedding, since the measured effective
dimension came out close to $N_1$ itself. We tested $1/\sqrt{N_1}$ and
$5/\sqrt{N_1}$ with 5-seed averaging: both performed \emph{worse}. We kept
$1/N_1$. This is the mirror image of the tabular finding --- there, the paper's
RKHS schedule failed to transfer to a tabular setting; here, the
tabular-calibrated schedule transferred fine to a genuine kernel.

\subsection{Full pessimism sweep}
\begin{center}\small
\begin{tabular}{@{}lccccc@{}}
\toprule
$c$ & 0.03 & 0.1 & 0.3 & 1.0 & 3.0 \\
\midrule
Fraction valid ($V_{\mathrm{low}} \le V_{\mathrm{true}}$) & 1.00 & 1.00 & 1.00 & 1.00 & 1.00 \\
Pessimistic regret & 0.479 & 0.438 & 0.397 & 0.397 & 0.397 \\
Plug-in regret & 0.241 & 0.241 & 0.241 & 0.241 & 0.241 \\
Max restart gap & 2.2e-16 & 1.36 & 5.26 & 20.4 & 62.2 \\
\bottomrule
\end{tabular}
\end{center}

Validity holds everywhere. Selection quality does not match the tabular case, and
the restart gap grows sharply with region width --- consistent with the
optimization landscape becoming genuinely harder to solve exactly, rather than
merely more conservative.

\section{Benchmark Sweep Figure}
\label{app:benchfig}

\begin{center}
\includegraphics[width=0.85\textwidth]{figures/fig6_mb_vs_mf.png}
\end{center}
\begin{center}
\small\textit{Naive vs.\ model-based (plain/shrunk) vs.\ model-free MAE, 5-seed
95\% CIs. Left: vs.\ confounding strength $\kappa$ at $N{=}20{,}000$. Right:
vs.\ sample size at $\kappa{=}1$. Naive stays lowest across both panels. The
wide, overlapping error bars on the two model-based variants are the visual
evidence behind the retracted shrinkage-variance claim of
\cref{app:corrections}.}
\end{center}

\section{Negative-Control Quality: Full Sweeps}
\label{app:nc}

Two metrics are reported at every grid point, because they disagree by
construction and reporting only one inverts the conclusion:
\emph{systematic bias} is the mean over policies of $|\,$mean over seeds of
$(\hat V - V)|$, isolating what the method is designed to remove; \emph{total
error} is the mean over policies and seeds of $|\hat V - V|$, which is what a
practitioner pays and includes the variance the correction costs. The plug-in
wins the first and loses the second. $N{=}2{,}000$, 5 seeds.

\begin{center}\small
\textbf{[N1] Strength: $\sigma_0$ varies, Assumption 3.1 holds throughout}\\[2pt]
\begin{tabular}{@{}cccccc@{}}
\toprule
& & \multicolumn{2}{c}{Systematic bias} & \multicolumn{2}{c}{Total error} \\
\cmidrule(lr){3-4}\cmidrule(lr){5-6}
$\sigma_0$ & Recov.\ err & plug-in & naive & plug-in & naive \\
\midrule
0.10 & 0.074 & \textbf{0.027} & 0.043 & 0.061 & 0.049 \\
0.25 & 0.083 & \textbf{0.029} & 0.043 & 0.067 & 0.049 \\
0.50 & 0.113 & \textbf{0.033} & 0.043 & 0.076 & 0.049 \\
1.00 & 0.191 & 0.046 & 0.043 & 0.084 & 0.049 \\
2.00 & 0.423 & 0.068 & 0.043 & 0.102 & 0.049 \\
\bottomrule
\end{tabular}
\end{center}

Naive is flat at $0.043$ throughout, as it must be --- it never touches $O_0$ ---
which makes it a clean control: everything that moves is attributable to the
negative control's quality.

\begin{center}\small
\textbf{[N2] Validity: $\delta$ breaks Assumption 3.1}\\[2pt]
\begin{tabular}{@{}cccccc@{}}
\toprule
& & \multicolumn{2}{c}{Systematic bias} & \multicolumn{2}{c}{Total error} \\
\cmidrule(lr){3-4}\cmidrule(lr){5-6}
$\delta$ & Recov.\ err & plug-in & naive & plug-in & naive \\
\midrule
0.00 & 0.083 & \textbf{0.029} & 0.043 & 0.067 & 0.049 \\
0.25 & 0.190 & \textbf{0.032} & 0.041 & 0.067 & 0.047 \\
0.50 & 0.378 & 0.039 & 0.038 & 0.068 & 0.047 \\
1.00 & 0.767 & 0.055 & 0.034 & 0.071 & 0.046 \\
\bottomrule
\end{tabular}
\end{center}

Caveats: $\delta$ implements one specific violation (leakage of the negative
control's own noise into the reward); others may behave differently. The
thresholds are properties of this environment, not universal constants --- the
transferable claim is the ordering, that validity matters more than strength.
Naive's bias also drifts slightly downward as $\delta$ grows
($0.043 \to 0.034$), so the widening gap is not purely the plug-in degrading.

\section{Corrections to Our Own Results}
\label{app:corrections}

We record these because the verification process producing them is part of the
contribution, not an embarrassment to be smoothed over.

\paragraph{Correction 1: harm ratio at zero confounding.} On a 3-seed sweep we
reported the model-based estimator as $23$--$46\times$ worse than naive at
$\kappa = 0$. At 5 seeds the ratio is $2.8$--$6.8\times$. Still a real, clear
negative result --- just a substantially smaller one than we had stated.

\paragraph{Correction 2: retracted variance-reduction claim.} We had claimed
James--Stein shrinkage ``cuts run-to-run variance sixfold under strong
confounding.'' It does not replicate. Across the full 15-cell sweep, shrinkage
has lower variance than the plain estimator in only $7/15$ cells --- roughly a
coin flip --- and at $\kappa = 1$ specifically, the exact condition cited, its
variance is \emph{higher} at $N{=}5{,}000$ and $N{=}20{,}000$. The claim is
retracted. The honest summary is a modest, noisy $\approx\!7\%$ average MAE
reduction, not a variance-reduction guarantee.

Neither was a coding bug: both were correctly computed standard deviations over
an $n{=}3$ sample, which is too few draws to be stable. The lesson we take
forward is that a claim derived from a \emph{standard deviation} needs more seeds
than a claim about a mean alone.

\paragraph{Correction 3: scope of the continuous study.} We initially described
the continuous extension as testing ``the paper's own generality.'' Re-reading
\citet{hong2024} showed it assumes continuous $S$ and $O$ but an explicitly
\emph{finite} action space, while our environment uses a real-valued action. The
oracle check still passes exactly, but the work is extrapolation beyond the
stated theorem rather than an instantiation of it. We corrected the framing
rather than relabel the experiment, and then ran the paper-faithful configuration
(\cref{sec:exp-fa}).

\paragraph{Correction 4: retracted mechanism for the selection failure.} We had
hypothesised that pessimistic selection underperformed because worst-case
perturbations compound through the $T{=}3$ chain, citing coordinate-descent
restart gaps that grew with region width. The finite-action experiment refutes
this on two counts: the anomaly persists when the action space is finite, and
restart gaps there are $\le 4\times10^{-16}$ up to $c{=}1$, so the optimizer is
finding the global minimum. The mechanism is instead behavior-policy coverage
(\cref{tab:coverage}). We report the failed hypothesis rather than silently
substituting the new explanation, because the sequence --- hypothesis, targeted
experiment, refutation, replacement --- is what makes the final claim credible.

\section{Continuous Extension: Per-Policy De-biasing}
\label{app:cont-tab}

Continuous kernel extension, 5 seeds, mean $\pm$ 95\% CI at $N{=}2{,}000$
(summarised in \cref{sec:exp-cont}). The pattern to note is in the intervals:
naive bias is tight but systematically wrong, plug-in bias is several times
smaller but noisier --- the bias-for-variance trade that recurs throughout this
paper.

\begin{center}\small
\begin{tabular}{@{}lcc@{}}
\toprule
Policy & Naive bias & Plug-in bias \\
\midrule
never\_treat & $-1.255 \pm 0.018$ & $\mathbf{-0.401 \pm 0.216}$ \\
constant\_dose & $+0.093 \pm 0.013$ & $\mathbf{+0.008 \pm 0.032}$ \\
proportional & $-0.435 \pm 0.013$ & $\mathbf{-0.129 \pm 0.061}$ \\
aggressive & $+1.357 \pm 0.038$ & $\mathbf{+0.365 \pm 0.178}$ \\
\bottomrule
\end{tabular}
\end{center}

\section{Finite-Action Experiment: Full Results}
\label{app:fa}

Environment: continuous scalar state and observation, $A = \{0, 0.5, 1.0\}$,
$T{=}3$, behavior policy a softmax over dose levels driven by the true hidden
state (so overlap holds by construction). Candidate policies are
observation-independent distributions over $A$; that restriction is what
preserves the exact oracle, since the expected bridge feature under such a policy
is exactly $[\pi, m_t]$ and the chain stays affine in the mean.

Two implementation improvements over \cref{app:cont} follow from the finite
action space. The bridge feature map $[\mathrm{onehot}(a), o]$ is
$(K{+}1)$-dimensional, so the primal solve is mandatory from the outset rather
than discovered. And because the chain is affine in the mean, dynamics-block
gradients become \emph{exact} via the backward recursion
$D_t = b_{R,t}[-1] + b_{D,t}[-1]\, D_{t+1}$, replacing the finite difference of
\cref{app:cont}; verified against numerical differentiation to
$1.16\times10^{-10}$.

\begin{center}\small
\begin{tabular}{@{}lcccc@{}}
\toprule
$N$ & 300 & 800 & 2{,}000 & 4{,}000 \\
\midrule
Reward-bridge error & 0.147 & 0.141 & 0.083 & \textbf{0.071} \\
Dynamics-bridge error & 0.271 & 0.251 & 0.156 & \textbf{0.086} \\
\bottomrule
\end{tabular}
\end{center}

Recovery is monotone in both families, cleaner than the continuous-action run,
which was non-monotone at an intermediate point.

\begin{center}\small
\begin{tabular}{@{}lccccc@{}}
\toprule
$c$ & 0.03 & 0.1 & 0.3 & 1.0 & 3.0 \\
\midrule
Fraction valid & 1.00 & 1.00 & 1.00 & 1.00 & 1.00 \\
Pessimistic regret & 0.211 & 0.211 & 0.289 & 0.602 & 0.602 \\
Plug-in regret & 0.000 & 0.000 & 0.000 & 0.000 & 0.000 \\
Max restart gap & 0.0 & 1.1e-16 & 2.2e-16 & 4.4e-16 & 5.33 \\
\bottomrule
\end{tabular}
\end{center}

Plug-in selection is perfect at every width; pessimistic selection is not, and
the near-zero restart gaps rule out an optimization explanation. See
\cref{tab:coverage} for the coverage diagnostic that does explain it.

\section{Bugs Found by Oracle Verification}
\label{app:bugs}

Two implementation bugs in the model-free comparator were invisible from the
derivation and surfaced only against exact ground truth. First, a per-stage
direct-method implementation was silently biased off-policy: a policy with true
value $1.865$ evaluated to $1.61$, with the error \emph{not} shrinking in $N$ ---
the fingerprint of a derivation error rather than noise. It was fixed with a
proper backward recursion. Second, the corrected version naively scaled to the
benchmark attempted to build and invert an $\approx\!11{,}500 \times 11{,}500$
dense matrix per policy per fit ($\approx\!2$\,GB), which stalled the run
entirely; exploiting a block-diagonal structure the naive implementation missed
resolved it.

The rank-deficiency issue of \cref{app:cont} belongs to the same category, with
the additional wrinkle that it was invisible even from the downstream recovery
metrics.

\section{Correction: the Pessimism Numbers on a Common Grid}
\label{app:pess-correction}

Three defects in our own implementation and two in our reporting were found after
the results of \cref{sec:exp-pess} were first written up. All are corrected here.

\paragraph{The omitted class bound.} \citet{hong2024} minimise over
$\mathrm{conf}(\alpha) \cap \mathcal{B}_{R,t}$, and the uniform bound $M_R$ on
that class carries through their Theorem~4.2. Our implementation minimised over
the confidence region alone. This is an error on our side, not in the theory, and
it inflates every pessimistic magnitude we reported.

\paragraph{The class carries two bounds, not one.} \citeauthor{hong2024}'s
Assumption~4.1(f) bounds the class in \emph{sup}-norm; their Assumption~D.16(d)
additionally norm-constrains it in
RKHS norm, which in this tabular instantiation is the coefficient $\ell_2$ norm.
Neither our ball nor our box alone is therefore the paper's literal class ---
each relaxes it, and the paper's class is their intersection, at least as strong
as the stronger of the two at every width. We report both. The box alone is the
weakest of the three: in this $36$-dimensional coefficient space the tightest box
containing the true bridge admits $\ell_2$ norms up to $5.2$ against the true
bridge's $1.803$.

\paragraph{Cross-$c$ reporting.} \cref{sec:exp-pess} quoted vanilla at $c{=}0.1$
against the projection at $c{=}0.03$. On a common grid the contrast disappears.

Toy POMDP, $T{=}3$, $N{=}20{,}000$, $3$ seeds, $V_{\mathrm{true}} = 2.0177$.
Selection regret on the left, $V_{\mathrm{low}}$ on the right:

\begin{center}
\begin{tabular}{lcccccccc}
\toprule
& \multicolumn{4}{c}{Selection regret} & \multicolumn{4}{c}{$V_{\mathrm{low}}$} \\
\cmidrule(lr){2-5}\cmidrule(lr){6-9}
$c$ & $0.03$ & $0.10$ & $0.30$ & $1.00$ & $0.03$ & $0.10$ & $0.30$ & $1.00$ \\
\midrule
Plug-in             & \textbf{0.000} & \textbf{0.000} & \textbf{0.000} & \textbf{0.000} & $2.052$ & $2.052$ & $2.052$ & $2.052$ \\
Vanilla             & $0.000$ & $0.257$ & $0.513$ & $0.513$ & $1.272$ & $0.123$ & $-3.235$ & $-20.983$ \\
$\ell_2$ ball       & $0.000$ & $0.257$ & $0.513$ & $0.513$ & $1.278$ & $0.469$ & $-0.806$ & $-3.026$ \\
Box (4.1(f) only)   & $0.000$ & $0.257$ & $0.427$ & $0.427$ & $1.275$ & $0.155$ & $-2.729$ & $-10.812$ \\
Signal-projected    & $0.000$ & \textbf{0.000} & $0.257$ & $0.513$ & $1.338$ & $0.813$ & $-0.298$ & $-4.015$ \\
\bottomrule
\end{tabular}
\end{center}

\paragraph{What this changes.} At $c{=}0.03$ --- the operating point at which
\cref{sec:exp-pess} reported the projection's $0.000$ --- \emph{every} method
including vanilla achieves $0.000$, so the published advantage is an artifact of
comparing across region sizes. The projection does have a real advantage in the
middle of the grid: at $c{=}0.1$ it alone among the pessimistic variants still
selects the optimal policy, and at $c{=}0.3$ its regret is half the others'.

\paragraph{What this does not change.} The ranking failure of vanilla pessimism
at larger $c$ survives the class correction, and the null-space mechanism of
\cref{sec:method-pess} is untouched. Identification, estimation and de-biasing
results do not involve the pessimism layer.

\paragraph{The baseline dominates.} The confounding-aware plug-in achieves
$0.000$ regret at every width tested and is never beaten by any pessimistic
variant. On this environment the pessimism layer costs selection accuracy and
returns nothing. We report this alongside the projection's $0.000$, because a
pessimistic method that does not beat its own plug-in has not earned its
complexity here.
